\section{Исследование литературы}

\subsection{Прогнозирование финансовых временных рядов}

Прогнозирование финансовых временных рядов является сложной задачей из-за шума, нестационарности, смены рыночных режимов и зависимости ценовой динамики от внешнего информационного фона. В классической финансовой теории эта сложность связана с тем, что цены быстро включают доступную рынку информацию \cite{fama1970efficient}. Поэтому задача прогноза не сводится только к выбору алгоритма: существенную роль играет то, какие именно признаки используются и какая информация считается доступной на момент прогноза.

Традиционные подходы к анализу финансовых рядов используют авторегрессионные и эконометрические модели, технические индикаторы, лаговые доходности, признаки волатильности и объёма торгов. Эти признаки описывают уже наблюдавшееся поведение рынка. С точки зрения машинного обучения такая постановка естественно связывается с задачей отбора и оценки признаков: новая группа признаков имеет смысл только тогда, когда она даёт дополнительный вклад относительно уже имеющихся переменных \cite{guyon2003feature}.

Позднее к классическим моделям добавились нейросетевые архитектуры. Для последовательных данных часто используется LSTM, поскольку эта архитектура предназначена для моделирования зависимостей между текущим и предыдущими состояниями временного ряда \cite{hochreiter1997lstm}. В финансовом прогнозировании LSTM может применяться как самостоятельная модель по ценовым данным или как компонент гибридной схемы, объединяющей рыночные и текстовые признаки.

Развитие Transformer-архитектур изменило подход к обработке последовательностей. В работе Vaswani et al. Transformer предложен как архитектура, основанная на attention-механизме и не требующая рекуррентной обработки \cite{vaswani2017attention}. Для данной работы это важно потому, что современные языковые модели, используемые для построения текстовых признаков, основаны именно на Transformer-family. BERT стал одной из ключевых моделей этого направления: он строит двунаправленные контекстные представления текста и может быть адаптирован к разным NLP-задачам \cite{devlin2019bert}.

\subsection{Текстовые данные как источник признаков}

Значительная часть информации, влияющей на финансовые рынки, существует в текстовой форме. К таким источникам относятся корпоративные раскрытия, финансовые новости, аналитические отчёты, сообщения инвесторов в социальных сетях и специализированные платформы вроде StockTwits. Текстовые данные могут отражать ожидания участников рынка, реакцию на новости, неопределённость, внимание к конкретному активу и интерпретацию событий.

Классические работы по финансовой текстовой аналитике показывают, что тексты могут быть связаны с рыночной динамикой. Tetlock показывает связь между медиа-содержанием и поведением фондового рынка \cite{tetlock2007media}. Loughran и McDonald показывают, что финансовый язык требует специальных словарей: слова, нейтральные в обычном языке, могут иметь другую интерпретацию в корпоративной отчётности \cite{loughran2011liability}. Обзоры Kearney и Liu, а также Nassirtoussi et al. фиксируют, что текстовая информация может использоваться как источник признаков, но результат зависит от источника текста, способа агрегации, временного выравнивания и целевой переменной \cite{kearney2014textual,nassirtoussi2014text}.

Ранние работы по социальным медиа также указывают на возможную связь между агрегированным настроением и рыночной динамикой. Bollen et al. анализируют Twitter-настроения и показывают, что они могут быть связаны с движением фондового индекса \cite{bollen2011twitter}. Mao et al. сравнивают несколько источников --- опросы, новости, Twitter и поисковые данные --- и тем самым показывают, что информационный фон может быть представлен разными текстовыми и поведенческими каналами \cite{mao2011sentiment}. Для данной работы эти результаты важны не как прямое доказательство прогнозируемости цены, а как обоснование того, что текстовый поток может быть преобразован в числовые признаки.

Ранние прикладные работы по корпоративным раскрытиям развивают ту же идею в более близкой к финансовой постановке. Feuerriegel и Fehrer рассматривают финансовые раскрытия компаний и показывают, что модели глубокого обучения могут использоваться для поддержки решений на основе таких текстов \cite{feuerriegel2015disclosures}. Kraus и Feuerriegel применяют transfer learning к финансовым раскрытиям \cite{kraus2017financialdisclosures}. Эти работы важны как предшественники современной схемы ``текст $\rightarrow$ признаки $\rightarrow$ прогноз''.

В более поздних исследованиях внимание смещается к пользовательским сообщениям и микроблогам инвесторов. Работы со StockTwits показывают, что пользовательский финансовый текст имеет собственную лексику, cashtag-структуру и специфический контекст \cite{jaggi2021stocktwits}. Поэтому для таких данных недостаточно универсальных инструментов анализа тональности: требуется либо доменная адаптация языковой модели, либо специальная проверка того, какие признаки действительно добавляют информацию к рыночным данным.

Для криптоактивов информационная среда может отличаться от фондового рынка. Криптовалюты торгуются непрерывно, без ограничений биржевой сессии; реакция на новости и социальный фон может быть быстрее, а роль narrative- и community-driven сигналов выше. Поэтому результат, полученный на акции отдельной компании, не должен автоматически переноситься на крипторынок. Проверка текстовых признаков на криптоактивах в данной работе рассматривается как отдельная проверка устойчивости вывода: она позволяет оценить, сохраняется ли эффект текстовых признаков при изменении информационной среды, источника текстов и горизонта прогноза.

\subsection{Языковые модели как генераторы признаков}

Переход от словарных методов и простых sentiment-score к контекстным языковым моделям стал важным этапом в финансовом NLP. BERT-подобные модели позволяют получать контекстные представления текста, которые учитывают не только отдельные слова, но и их значение в конкретной фразе. Это особенно важно для финансового языка, где одни и те же слова могут иметь разную интерпретацию в зависимости от контекста.

Для финансовых задач особое значение имеет FinBERT \cite{araci2019finbert}. Эта модель является адаптацией BERT к финансовым текстам и используется для задач финансовой тональности и построения текстовых представлений. Дальнейшая линия работ по FinBERT для финансовых коммуникаций также показывает, что доменная адаптация BERT может улучшать качество анализа финансовых текстов относительно универсальной языковой модели \cite{yang2020finbert}.

Важно различать несколько способов использования FinBERT. Первый --- извлечение sentiment-классов или sentiment-score: positive, negative, neutral. Второй --- использование скрытых представлений модели, то есть embeddings. Третий --- построение более простых text-meta признаков: длина текста, число сообщений, интенсивность новостного потока и другие агрегированные характеристики. В первом случае текст сводится к компактной оценке тональности. Во втором случае текст представлен многомерным вектором, который может содержать более широкий набор сигналов: контекст, упоминания событий, неопределённость, ожидания и семантические связи. В третьем случае проверяется не смысл текста напрямую, а интенсивность и структура текстового потока.

Для данной работы это различие принципиально. В основном эксперименте и в дополнительной проверке с расширенным рыночным baseline используются FinBERT embeddings как высокоразмерные текстовые представления. В мультиактивной crypto-проверке используются более компактные FinBERT-derived sentiment-признаки и text-meta признаки. Такое различие не является противоречием: работа исследует не одну конкретную архитектуру, а класс LLM/FinBERT-generated features и их добавочный вклад относительно рыночных признаков.

\subsection{Прикладные мультимодальные исследования}

Прикладная литература показывает, что наиболее содержательной является не чисто текстовая, а мультимодальная постановка, в которой текстовые признаки сопоставляются с рыночными baseline. В работе Hiew et al. BERT-based sentiment index используется вместе с другими каналами настроений и рыночными факторами в LSTM-модели \cite{hiew2019bert}. Такая схема близка к логике настоящей работы: текст сначала преобразуется в числовые признаки, а затем проверяется как возможное дополнение к рыночной информации.

Отдельно важна линия работ, где текстовые признаки объединяются с ценами, техническими индикаторами и другими рыночными переменными. Jiang и Zeng используют FinBERT для извлечения сентимента из финансовых новостей и объединяют эти признаки с рыночными данными в LSTM-постановке \cite{jiang2023finbert}. Однако подобные результаты необходимо интерпретировать осторожно: регрессия цены закрытия, классификация направления и классификация тональности являются разными задачами, поэтому их метрики нельзя напрямую сравнивать.

Наиболее близким к RUN 05 является исследование PreBit, где для Bitcoin объединяются OHLCV-данные, технические индикаторы, связанные активы и Twitter FinBERT embeddings \cite{zou2022prebit}. Важная особенность этой работы состоит не только в использовании Bitcoin, но и в явной мультимодальной логике: авторы проверяют вклад разных модальностей через ablation study. Поэтому PreBit является полезным ориентиром для данной курсовой: он показывает, что для криптоактивов текстовые признаки следует оценивать не изолированно, а относительно рыночной модели и других групп признаков.

Современные работы также показывают, что увеличение сложности языковой модели само по себе не гарантирует улучшения прогноза. Shobayo et al. сравнивают FinBERT, GPT-4-подход и логистическую регрессию на финансовых новостях и показывают, что простая модель может быть конкурентоспособной или даже сильнее сложной языковой схемы \cite{shobayo2024sentiment}. Albada и Sonola используют BERTweet и Transformer-подходы для задачи next-day stock movement на StockNet, но при честном временном разбиении абсолютные улучшения остаются умеренными \cite{albada2025bertweet}. Vuong и Gauch используют LLM не как прямой предсказатель цены, а как инструмент улучшения tweet-emotion признаков перед downstream-моделью \cite{vuong2025llm}. Эти работы поддерживают feature-centric подход: ценность языковой модели определяется тем, даёт ли она переносимый признак, а не числом параметров самой модели.

Следовательно, исследовательский пробел состоит не в отсутствии ещё одной сложной архитектуры для прогноза цены. Более существенная задача --- проверить добавочную предсказательную силу разных групп признаков: рыночных признаков, embeddings, sentiment-признаков и text-meta признаков. Именно такую постановку использует данная работа: модели применяются как инструмент сравнения признаков, а не как самоцель.

\subsection{Методологические ограничения существующих исследований и место данной работы}

Обзор литературы показывает, что исследования по текстовому прогнозированию финансовых временных рядов часто отличаются не только моделями, но и самой постановкой задачи. Одни работы прогнозируют цену или доходность, другие --- бинарное направление движения, третьи --- экстремальные движения, четвёртые --- только тональность текста. Поэтому прямое сравнение метрик между статьями ограничено: accuracy, MSE, F1-score и ROC-AUC могут относиться к разным задачам.

Первая методологическая проблема связана с разбиением данных. В задачах временных рядов случайное разделение наблюдений может привести к утечке информации, потому что обучающая и тестовая подвыборки будут содержать близкие по времени наблюдения. Более надёжным является out-of-time разбиение, при котором модель обучается на раннем периоде и тестируется на более позднем. В данной работе используется хронологическое разбиение без пересечения дат между обучающей, валидационной и тестовой подвыборками.

Вторая проблема --- предварительная обработка данных. Если scaler, PCA или другие преобразования обучаются на всей выборке, в модель попадает информация о распределении будущих данных. Поэтому в корректном протоколе такие преобразования должны обучаться только на обучающей подвыборке, а затем применяться к валидационной и тестовой подвыборкам без повторной подгонки. В данной работе это условие явно фиксируется в экспериментальном протоколе.

Третья проблема --- отсутствие сильных baseline. Ряд исследований показывает результат сложной модели, но не всегда сопоставляет его с majority baseline, рыночной моделью или простой линейной моделью. В результате трудно понять, действительно ли текстовые признаки дают добавочную предсказательную силу. В данной работе эта проблема решается через прямое сравнение групп признаков: рыночные признаки, FinBERT/текстовые признаки и их объединение.

Четвёртая проблема связана с множественным сравнением. Если исследователь одновременно перебирает несколько активов, горизонтов, моделей и групп признаков, максимальный найденный эффект может быть завышен. В такой ситуации положительный результат по одному активу или одному горизонту нельзя автоматически трактовать как универсальное доказательство полезности текстовых признаков. Такие результаты требуют проверки на независимом периоде, через rolling-window validation или через специальные процедуры контроля множественных сравнений.

На этом фоне место данной работы можно определить как feature-centric эксперимент, проверяющий добавочную силу LLM/FinBERT-generated features относительно рыночных признаков. Основной эксперимент на AAPL показывает базовую постановку сравнения, RUN 04 проверяет устойчивость вывода при усилении рыночного baseline, а RUN 05 расширяет анализ на crypto-news и несколько горизонтов прогноза. Работа не претендует на универсальный вывод о всех рынках и всех языковых моделях; её задача состоит в том, чтобы показать, при каких условиях текстовые признаки дают или не дают дополнительный прогнозный вклад.
